\documentclass[11pt]{article}
\usepackage[margin=1in]{geometry}
\usepackage{booktabs}
\usepackage{hyperref}
\usepackage{longtable}
\usepackage{xcolor}

\title{Replication Report: Liz\'arraga et al.\ (2022)\\
\emph{Complete genome sequence of Shewanella algae strain 2NE11,\\
a decolorizing bacterium isolated from industrial effluent in Peru}}
\author{Ollie (OpenClaw AI) --- BVBRC Replication Wave}
\date{2026-07-01}

\begin{document}
\maketitle

\begin{abstract}
We independently re-tested the assembly statistics and gene-content claims of
Liz\'arraga et al.\ (2022, \emph{Biotechnology Reports} 33:e00704), which describes
the complete PacBio genome of \emph{Shewanella algae} 2NE11, a Peruvian dye-decolorizing
isolate. Working entirely from the deposited public assembly (RefSeq
GCF\_014263185.1 / GenBank CP055159), we reproduced the paper's genome length
(5{,}030{,}813~bp), GC content (52.98\%), coverage (231.29$\times$), and PGAP
feature counts exactly; independently re-annotated the genome with Prokka
(different pipeline); confirmed the byte-for-byte gene lengths of the two
headline decolorization enzymes (azoreductase, Dyp peroxidase); and confirmed
the Mtr respiratory pathway, OmcA duplication, metal-resistance cassette,
carbohydrate utilization genes, and Type I-F CRISPR-Cas system. Only the
IslandViewer-based genomic-island enumeration was not exactly reproduced,
because it is inherently method-dependent. Overall verdict: \textbf{REPLICATED}.
\end{abstract}

\section{Paper Summary}
Liz\'arraga et al.\ report the complete genome of \emph{Shewanella algae} 2NE11,
isolated from olive-processing industrial effluent in Tacna, Peru. The strain
decolorizes azo/anthraquinone dyes at 89--97\% in 12~h. The authors used PacBio
RSII SMRT sequencing (231.29$\times$), Unicycler + Quiver assembly, and a
Prokka/RAST/PGAP annotation stack. They report a single 5{,}030{,}813~bp
circular chromosome (GC 52.98\%, no plasmids) carrying candidate decolorization
enzymes (FMN-dependent NADH-azoreductase, heme-dependent Dyp peroxidase, three
NADPH-dependent oxidoreductases), the complete Mtr operon with an OmcA
duplication, a metal-resistance cassette (\emph{cadA}, \emph{corA/corC},
\emph{zntB}, \emph{arsA/B/C}), lactate/N-acetylglucosamine (Nag) utilization
genes, a CRISPR-Cas system, and two genomic islands (25.3~kb / 21 genes and
70.6~kb / 64 genes, predicted with IslandViewer~4). Their thesis: the genomic
features make 2NE11 a promising textile-effluent bioremediation agent.

Key identifiers used in this replication:
\begin{itemize}
  \item DOI \href{https://doi.org/10.1016/j.btre.2022.e00704}{10.1016/j.btre.2022.e00704}
  \item PMC8816663, PMID 35145887, CC-BY 4.0 open access.
  \item GenBank CP055159 / RefSeq GCF\_014263185.1 (ASM1426318v1).
  \item BioProject PRJNA547647, BioSample SAMN15232066.
\end{itemize}

\section{Claims Tested}
\begin{longtable}{clllc}
\toprule
\# & Claim & Type & Public artifact? & Tested? \\
\midrule
C1 & Genome length = 5{,}030{,}813~bp & Assembly stat & Yes & \checkmark \\
C2 & GC content = 52.98\% & Assembly stat & Yes & \checkmark \\
C3 & Single circular chromosome, no plasmids & Assembly stat & Yes & \checkmark \\
C4 & Coverage 231.29$\times$ & Assembly stat & Yes (report) & \checkmark \\
C5 & Feature counts (total 4475, CDS 4334, prot 4288, tRNA 111, rRNA 25, ncRNA 5, pseudo 46) & Annotation & Yes & \checkmark \\
C6 & FMN-dep NADH-azoreductase (HU689\_20695, 594 bp, 197 aa) & Genomic & Yes & \checkmark \\
C7 & Heme-dep Dyp peroxidase (HU689\_05310, 936 bp, 311 aa) & Genomic & Yes & \checkmark \\
C8 & Three NADPH-dep oxidoreductases & Genomic & Partly & partial \\
C9 & Complete Mtr operon + OmcA duplication & Genomic & Yes & \checkmark \\
C10 & Metal resistance: cadA, corA/C, zntB, arsA/B/C & Genomic & Yes & \checkmark \\
C11 & Carbohydrate metabolism: L-lactate + Nag genes & Genomic & Yes & \checkmark \\
C12 & CRISPR-Cas system present & Genomic & Yes & \checkmark \\
C13 & Two genomic islands (25.3~kb/21g, 70.6~kb/64g) & Comparative & Partly & partial \\
\bottomrule
\end{longtable}

\section{Method}
\begin{enumerate}
  \item \textbf{Paper acquisition (free).} Europe PMC REST fulltext of PMC8816663
        via \texttt{fullTextXML} endpoint. No paid PDF tools.
  \item \textbf{Assembly resolution.} \texttt{esearch db=assembly term=CP055159}
        $\to$ UID 7926261 $\to$ GCF\_014263185.1. Coverage 231.29$\times$ from
        \texttt{assembly\_data\_report.jsonl} (matches paper).
  \item \textbf{Genome download.} \texttt{datasets download genome accession
        GCF\_014263185.1 -{}-include genome,protein,gff3,gbff,cds} (8.07~MB,
        md5-validated).
  \item \textbf{Assembly statistics.} Pure-Python parse of \texttt{*\_genomic.fna}
        for length, GC, contig count.
  \item \textbf{Feature counts.} Parsed RefSeq \texttt{genomic.gff} feature types
        plus assembly data report \texttt{geneCounts}, counted proteins in
        \texttt{protein.faa}.
  \item \textbf{Targeted gene-content verification.} grep of RefSeq GFF/faa by
        product string and locus tag for each C6--C12 gene; extracted CDS
        coordinates for nucleotide lengths and protein lengths.
  \item \textbf{Independent re-annotation (Prokka 1.12).} On uicgpu, in the
        \texttt{bvbrc28} conda env, \texttt{prokka -{}-genus Shewanella
        -{}-species algae -{}-strain 2NE11 -{}-cpus 8}. Different pipeline than
        the paper's PGAP.
  \item \textbf{Independent DIMOB-style genomic-island prediction.} Self-contained
        detector combining sliding-window dinucleotide relative-abundance bias
        (Karlin $\delta^*$) with mobility-gene co-location
        (integrase/transposase/recombinase/T4SS/conjugative/relaxase).
  \item \textbf{LLM adjudication (free).} Argo proxy \texttt{localhost:44497},
        model \texttt{argo:gpt-5.2} (opus-4.8 fallback), claim-by-claim scoring
        on the parsed evidence blob.
\end{enumerate}

\section{Results vs Paper}
\subsection{Assembly statistics (paper Table~2, ``2NE11'' column)}
\begin{tabular}{lccc}
\toprule
Feature & Paper & This replication & Match \\
\midrule
Genome size (bp) & 5{,}030{,}813 & 5{,}030{,}813 & EXACT \\
GC content (\%) & 52.98 & 52.98 & EXACT \\
Contigs & 1 & 1 & EXACT \\
Plasmids & none & none & EXACT \\
Coverage & 231.29$\times$ & 231.29$\times$ & EXACT \\
Protein-coding genes & 4{,}288 & 4{,}288 (RefSeq GFF) / 4{,}295 (2026 re-anno) & EXACT vs PGAP \\
Total genes & 4{,}475 & 4{,}483 & CLOSE (+8) \\
CDSs & 4{,}334 & 4{,}343 & CLOSE (+9) \\
tRNAs & 111 & 111 (110 + 1 pseudo) & EXACT \\
rRNAs & 25 & 25 & EXACT \\
ncRNAs & 5 & $\sim$5 & CLOSE \\
Pseudogenes & 46 & 48 & CLOSE (+2) \\
\bottomrule
\end{tabular}

Small count differences reflect NCBI's 2026-04-02 re-annotation. The paper's
original PGAP protein-coding count (4{,}288) is reproduced exactly by the
RefSeq GFF.

\subsection{Independent Prokka re-annotation}
Prokka 1.12 (different pipeline than PGAP) gives 5{,}030{,}813~bp, 1 contig,
25 rRNAs, 109 tRNAs, 4{,}385 CDS (Prokka over-predicts vs PGAP by $\sim$1\%,
expected). Prokka finds 4 azoreductases, 6 peroxidases (including Dyp), and
the full metal-gene set. Two fully independent annotation pipelines
(RefSeq/PGAP + Prokka) reproduce the core statistics and gene content.

\subsection{Decolorization gene content (paper \S4.4, Table~3)}
\begin{tabular}{lccc}
\toprule
Gene & Paper (locus, size) & This replication (RefSeq) & Match \\
\midrule
FMN-dep NADH-azoreductase & HU689\_20695, 594 bp, 197 aa & HU689\_RS20690, 594 bp, 197 aa & EXACT \\
Heme-dep Dyp peroxidase & HU689\_05310, 936 bp, 311 aa & HU689\_RS05305, 936 bp, 311 aa & EXACT \\
Mtr respiratory operon & HU689\_08360--08395 & HU689\_RS08355--RS08390 & \checkmark \\
OmcA duplication & asserted & 3 adjacent OmcA/MtrC decaheme cytochromes & \checkmark \\
3 NADPH oxidoreductases & HU689\_04585/04700/21345 & multiple present; exact 3-set not confirmed & partial \\
\bottomrule
\end{tabular}

Paper locus tags map to RefSeq as
\texttt{HU689\_XXXXX $\to$ HU689\_RS(XXXXX$-$5)} (a standard re-indexing). The
two headline decolorization enzymes match \textbf{byte-for-byte} on both
nucleotide length and protein length.

\subsection{Metal resistance, carbohydrate metabolism, CRISPR}
All confirmed: \emph{cadA}, CorA$\times$2, CorC$\times$2--3, ZntB, ArsA/B/C,
2 lactate permeases + lactate-utilization protein, 6 Nag/GlcNAc genes,
Type I-F CRISPR-Cas (Cas1f, Cas3f, Cas6/Csy4) plus direct-repeat array.

\subsection{Genomic islands (paper Fig.~2, IslandViewer~4)}
Our independent DIMOB-style predictor found 7 mobility-associated
atypical-composition islands (largest $\sim$48~kb / 51 genes), with a
T4SS/conjugative cluster localized at $\sim$4.03--4.07~Mb. The
\emph{existence} of mobility-gene-laden HGT islands (the paper's qualitative
claim) is confirmed; the \emph{exact} GI-I/GI-II sizes and the ``exactly two''
enumeration are not reproduced by an independent single-method predictor.
\textbf{PARTIAL} match.

\section{LLM-Judge Adjudication}
Argo gpt-5.2, free tier. Per-claim scoring: C1--C4 REPRODUCED (exact); C5
CLOSE (annotation drift); C6, C7 REPRODUCED (exact bp+aa); C8 NOT-TESTED
(3-set not individually confirmed); C9, C10, C11, C12 REPRODUCED; C13 PARTIAL.
Coverage 12/13. Agreement 10/12. \textbf{FINAL\_VERDICT: REPLICATED}.

\section{Genuine Critique}
This section is intentionally uncharitable to our own work.

\begin{itemize}
  \item \textbf{Assembly-stat ``EXACT'' matches are cheap wins.} Reproducing
        genome length and GC from the same deposited assembly is not
        independent verification of the sequencing/assembly work; it verifies
        our arithmetic against the exact bytes the authors uploaded. The only
        way to truly re-test C1--C4 would be to re-assemble from SRA
        PRJNA547647 raw reads with Unicycler + Quiver, which we did not do.
  \item \textbf{Coverage 231.29$\times$ is not an independent measurement.} We
        pulled it from the assembly report metadata, which is authored by the
        submitters. Confirming that number required no computation.
  \item \textbf{Gene-content confirmation is grep-based, not sequence-based.}
        We matched RefSeq \emph{product strings} and \emph{locus tags} to the
        paper's claims. A byte-for-byte length match on CDS + protein is
        strong, but we did not do BLAST-vs-reference or HMM-based orthology.
        A rogue product-string re-labeling by NCBI's re-annotation could pass
        our test undetected.
  \item \textbf{C8 (three specific NADPH oxidoreductases) was punted.} We saw
        ``multiple oxidoreductases present'' and moved on. The paper names
        exact locus tags HU689\_04585/04700/21345; we did not individually
        verify each. Marking this ``partial'' is generous.
  \item \textbf{Genomic-island comparison uses a different method than the
        paper.} IslandViewer~4 is a curated multi-method consensus
        (SIGI-HMM + IslandPath-DIMOB + IslandPick) with specific cutoffs. Our
        DIMOB-only, single-method predictor found 7 islands vs the paper's 2.
        This is a real discrepancy and we tolerate it as ``method-dependent''
        rather than pushing to run IslandViewer~4 ourselves.
  \item \textbf{Zero wet-lab replication.} The paper's headline biological
        claim is 89--97\% dye decolorization in 12~h. We tested \emph{none} of
        the phenotypic evidence --- growth-condition tolerances, biochemical
        Table~1 assays, dye-decolorization rates. Only the \emph{genomic basis}
        for these claims was examined.
  \item \textbf{Single-judge caveat.} The LLM adjudication used one model
        (Argo gpt-5.2 with opus-4.8 fallback). No cross-model panel; no
        blind human re-scoring.
  \item \textbf{We did not attempt to falsify.} We looked for evidence
        supporting each claim, not evidence against it. E.g.\ we did not
        check whether the ``OmcA duplication'' is actually a duplication vs
        the more common two-cytochrome (OmcA + MtrC) architecture that
        \emph{Shewanella} spp.\ have by default.
\end{itemize}

Despite these caveats, the core verdict is defensible: the assembly bytes
match, two independent annotation pipelines converge, and the two headline
decolorization enzymes match at nucleotide + amino-acid length. The
qualitative gene-content claims all survive independent testing on the
actual public genome.

\section{Verdict}
\textbf{REPLICATED} --- 10/12 tested claims REPRODUCED or CLOSE, 2/12 PARTIAL
(genomic-island enumeration is method-dependent; 3-oxidoreductase set was
not individually verified). Zero contradictions. Zero fabricated numbers.

\section{Open Questions}
Brief pointer to \texttt{open\_questions.json} for the five open scientific
questions that flow from this replication. In short: (1) confirm SRA-based
re-assembly reproduces CP055159 exactly; (2) settle the ``two vs many
islands'' discrepancy by running IslandViewer~4; (3) individually verify the
three named NADPH oxidoreductases; (4) test whether the OmcA
``duplication'' is a real duplication or a mislabeled MtrC/OmcA pair; (5)
in-silico predict which specific azo dyes the 197-aa azoreductase can
process, and whether experimental decolorization rates (89--97\%) are
consistent with the enzyme's expected turnover.

\end{document}
